\subsection{Data}\label{sec:data}

We evaluate the method on three Alpine glaciers that span a wide range of size and
are among the most intensively surveyed in the world (Table~\ref{tab:data}): Great
Aletsch (Swiss Alps), the largest glacier of the Alps ($\approx$80~km$^2$); Rhône
(Swiss Alps), a mid-sized valley glacier ($\approx$15~km$^2$) that feeds the
eponymous river; and Argentière (French Alps, Mont-Blanc massif), a smaller, steep
and fast-flowing valley glacier ($\approx$12~km$^2$). Each combines repeat DEMs,
satellite surface-velocity fields and a decades-long in situ stake
network---maintained by swisstopo and GLAMOS in Switzerland and by GLACIOCLIM in
France---which supply both the fields that drive the inversion and the ground truth
that validates it.

For each glacier the pipeline requires two surface DEMs bracketing the inversion
window---whose difference is the geometry change to be explained---and a
surface-velocity field for the data assimilation (Sect.~\ref{sec:da}).
Ice-thickness soundings are used only where available, as an optional refinement
(Appendix~A1). For Argentière the two DEMs are the 2012 and 2021 Pléiades
models of \citet{kneib2024}; for Great Aletsch (2009, 2017) and Rhône (2007, 2016)
they are swisstopo geodetic DEMs. We resample all fields onto a common
$100\,\mathrm{m}$ grid. Surface velocities are the satellite feature-tracking
mosaics of \citet{millan2022}, built from Sentinel-2 and Landsat-8 image pairs over
2017--2018. Because alpine surface velocities vary slowly relative to the decadal DEM
windows \citep{dehecq2019}, this single field serves as a representative snapshot for
the thickness calibration.

For validation we use in situ stake balances, kept strictly independent of the
inversion. Rather than compare against every raw stake reading---noisy and
irregularly sampled in space and time---we retain only locations with a continuous
multi-year record, average each into a single balance value ($\pm$ interannual
standard deviation), and evaluate the inferred smb against these averages.
This yields two such locations on Aletsch and five on Rhône; on Argentière the denser GLACIOCLIM
network provides the published 2012--2021 mean balance at each of 25 stakes.

\begin{table}
  \caption{Input and validation data for the three study glaciers. DEM epochs
    bracket the inversion window and velocities drive the data assimilation. The
    stake column gives the number of continuous-record locations averaged for the
    sRMSE.}
  \label{tab:data}
  \centering
  \begin{tabular}{lccc}
    \toprule
     & Argentière & Great Aletsch & Rhône \\
    \midrule
    Window            & 2012--2021 & 2009--2017 & 2007--2016 \\
    Surface DEMs      & Pléiades & swisstopo & swisstopo \\
    Velocity          & \multicolumn{3}{c}{\citep{millan2022}} \\
    Stake locations   & 25         & 2         & 5          \\
    \bottomrule
  \end{tabular}
\end{table}
